% PROFESSOR.tex --- Os cantos do homem-sombra (1ª série do Ensino Médio)
% Orientações para o professor (entra apenas no Livro do Professor)
% (TR SEDUC-GO 118787/2026, itens 6.4.2 e 6.4.3)

\blankAtodd
\part{Orientações para o professor}

\chapter{Carta ao professor}

Professora, professor,

\textit{Os cantos do homem-sombra} abre um percurso de leitura de três
anos dedicado a narrativas de povos originários do Brasil. Na 1ª série,
a turma conhece uma narrativa breve dos Hupd'äh, do Alto Rio Negro,
sobre a origem dos cantos de festa. Na 2ª série, lê \textit{A árvore dos
cantos}, um ciclo de narrativas de transformação contadas por pajés
yanomami. Na 3ª série, chega a \textit{A terra uma só}, em que o Guarani
Mbya Timóteo Verá Tupã Popygua escreve diretamente em português a
cosmogonia e a história de seu povo. A progressão vai, assim, do canto
breve e da oralidade coletiva à escrita autoral indígena, e da
sensibilização à análise crítica.

A brevidade deste livro é uma vantagem pedagógica: permite leituras
repetidas, em voz alta, e um trabalho minucioso com a linguagem, com a
tradução e com o livro como objeto bilíngue. A partir dele, é possível
discutir com estudantes de 15 ou 16 anos questões que atravessam todo o
Ensino Médio: o que é conhecimento, quem tem autoridade para contar uma
história, por que tantas línguas estão ameaçadas, como se formam e se
desfazem os estereótipos sobre os povos indígenas.

As atividades do estudante estão organizadas em etapas (antes, durante e
depois da leitura; aprofundamento; criação; pesquisa e debate; conexões
com Goiás; projeto integrador). Estas orientações oferecem subsídios
teóricos, um planejamento aula a aula, respostas esperadas e sugestões
de avaliação. Adapte-as livremente à realidade da sua turma.

\chapter{Sobre a obra}

\section{o povo hupd'äh}

Os Hupd'äh vivem na região do Alto Rio Negro, no estado do Amazonas, na
fronteira com a Colômbia. Segundo a nota dos organizadores, são cerca de
1\,500 pessoas distribuídas em 35 aldeias. Diferentemente de povos que
se fixaram às margens dos grandes rios, os Hupd'äh ocupam
tradicionalmente o interior da floresta, junto aos igarapés, e se
deslocam por longos caminhos para visitar parentes. Caçam com arco e
flecha e zarabatana, pescam nos igarapés e cultivam roças de maniva.

A língua hup, primeira língua das crianças, pertence à família naduhup,
com as línguas dos Yuhupdëh, dos Dâw e dos Nadëb. Esses povos foram por
muito tempo designados, na literatura e na região, como ``Maku'', termo
dado por povos vizinhos e hoje considerado pejorativo. Vale a pena
explorar com a turma a relação entre nomeação e hierarquia: no Alto Rio
Negro, as relações entre os povos do rio e os povos da floresta foram
historicamente assimétricas, e os nomes carregam essa história.

Desde 2001, professores hupd'äh participam de cursos de formação para
suas escolas, nos quais, com linguistas e antropólogos, descreveram a
língua e criaram uma escrita para ela. Foram elaborados um dicionário
hupd'äh--português e cartilhas.

\section{como foi feito este livro}

Em 2002, a linguista Patience Epps, então morando na aldeia de Barreira
Alta, gravou a narração do senhor Mário Andrade Pires. O texto foi
transcrito e traduzido com a ajuda dos professores hupd'äh e integrou o
livro \textit{Hupd'äh nɨh Pɨnɨ̀gd'äh} (\textit{Contos dos Hupd'äh}), usado
na alfabetização em língua hup. O canto \textit{Way Naku} foi cantado em
2010, na aldeia de Tat Dëh, pelo senhor Ponciano Salustiano Ramos para o
antropólogo Danilo Paiva Ramos, que coorganiza o volume. As ilustrações
são de Anita Ekman.

Esse percurso é, em si, conteúdo de aula. Ele mostra que a obra tem
\textit{autoria coletiva}: pertence a uma tradição transmitida por
gerações, foi narrada por anciãos identificados pelo nome, registrada e
traduzida por pesquisadores em parceria com professores indígenas, e
editada para dois públicos: as escolas hupd'äh e os leitores não
indígenas. O caráter bilíngue da edição, com os dois textos na mesma
página, afirma a igualdade de estatuto entre as línguas.

\chapter{Gênero, linguagem e temas}

\vspace*{-2\baselineskip}

\section{uma narrativa de tradição oral}

O texto pertence ao campo das narrativas de tradição oral. É
também uma \textit{narrativa de origem}: explica como os Hupd'äh
passaram a conhecer os \textit{caapivaiá}, os cantos de festa. Evite
apresentá-lo como ``lenda'', ``fábula'' ou ``folclore'', termos que, no
uso escolar, costumam sugerir fantasia sem valor de verdade ou curiosidade
exótica. Para os Hupd'äh, trata-se de conhecimento: um saber sobre a
origem de práticas que continuam vivas e sobre como se comportar diante
das outras gentes que habitam a mata.

O antropólogo Eduardo Viveiros de Castro observa que o interesse
ocidental pelos ``conhecimentos tradicionais'' costuma buscar apenas seus
conteúdos, sem se deixar transformar por sua forma, isto é, pela própria
ideia de conhecimento, de quem conhece e de para que se conhece. Vale
propor à turma essa inversão: não perguntar só ``o que essa história
diz?'', mas ``que tipo de saber é esse, como ele se transmite e o que ele
nos obriga a repensar?''.

\section{marcas de oralidade}

Mesmo traduzido e escrito, o texto preserva marcas da oralidade que
merecem ser observadas com os estudantes:

\begin{itemize}
\item \textbf{Fórmulas de transmissão}: ``Dizem que o Way
Naku ficou cantando\ldots'' e ``Foi isso o que ouvi de meus avós''.
Essas fórmulas situam a história numa cadeia de narradores e lhe
conferem autoridade.
\item \textbf{Alternância de vozes}: a narração em terceira
pessoa dá lugar à primeira pessoa do plural (``nós, Hupd'äh'', ``nossas
festas''), revelando um narrador que fala de dentro da comunidade.
\item \textbf{Onomatopeias}: ``Tac!'', ``Puffff!'',
``Vrummm!'', ``tchibum!'', que no texto hup aparecem como \textit{Täk!},
\textit{Pë'!}, \textit{Mmmm'!} e \textit{tapuh!}. Sua comparação é uma
porta de entrada concreta para discutir tradução.
\item \textbf{Frases curtas, repetições e economia
narrativa}: não há descrições psicológicas extensas; as motivações
ficam implícitas (por que o homem hup bate na árvore?), abrindo espaço
para a interpretação.
\item \textbf{Cantos com refrão}: os dois cantos de abertura
repetem versos e mantêm, mesmo na versão em português, um refrão em hup
(\textit{Way Naku, yari nóóy mah / Marika, Way Naku}), sinal de que
certas palavras não se traduzem sem perda.
\end{itemize}

\section{temas}

\begin{itemize}
\item A origem dos cantos e festas e o papel da música na vida coletiva.
\item A aprendizagem pela escuta atenta e pela imitação (``Escutou e
prendeu os cantos'').
\item As outras gentes da mata e a ética do encontro: perigo, fascínio,
respeito e astúcia.
\item A transmissão intergeracional do conhecimento.
\item Oralidade, escrita, tradução e línguas indígenas.
\item Diversidade dos povos indígenas e combate a estereótipos.
\end{itemize}

\chapter{Objetivos, competências e habilidades}

\section{objetivos}

\begin{itemize}
\item Ler, compreender e interpretar uma narrativa de tradição oral
hupd'äh, reconhecendo seus recursos expressivos e seu contexto de
produção.
\item Reconhecer a autoria coletiva e o processo de registro, tradução e
edição de narrativas orais.
\item Compreender a diversidade étnica e linguística dos povos indígenas
no Brasil, em especial no Alto Rio Negro e em Goiás.
\item Identificar e questionar estereótipos e preconceitos sobre os
povos indígenas.
\item Produzir textos, imagens e apresentações que dialoguem de modo
criativo, crítico e ético com a obra.
\end{itemize}

\section{fundamentos legais e curriculares}

O trabalho atende à Lei 11.645/2008, que torna obrigatório o estudo da
história e da cultura afro-brasileira e indígena em todo o currículo, e
às Diretrizes Curriculares Nacionais para a Educação das Relações
Étnico-Raciais. Dialoga, ainda, com o Documento Curricular para Goiás --
Etapa Ensino Médio (\textsc{dc-goem}), que orienta a abordagem da diversidade
étnico-racial e das culturas locais nas áreas do conhecimento.

\textbf{Competências gerais da BNCC mobilizadas}: 1
(conhecimento), 3 (repertório cultural), 4 (comunicação), 7
(argumentação) e 9 (empatia e cooperação).

\textbf{Temas Contemporâneos Transversais}: Multiculturalismo
(diversidade cultural e valorização das matrizes históricas e culturais
brasileiras), Meio Ambiente (educação ambiental) e Cidadania e Civismo
(educação em direitos humanos).

\section{habilidades da bncc}

\textbf{Linguagens e suas Tecnologias}

\begin{itemize}
\item \textbf{\textsc{em13lp01}}: relacionar o texto às suas condições de
produção e ao seu contexto sócio-histórico de circulação.
\item \textbf{\textsc{em13lp06}}: analisar efeitos de sentido de usos expressivos
da linguagem (onomatopeias, repetições, alternância de vozes).
\item \textbf{\textsc{em13lp10}}: analisar a variação linguística e fundamentar o
respeito à diversidade de línguas e variedades, combatendo preconceitos
linguísticos.
\item \textbf{\textsc{em13lp30}} e \textbf{\textsc{em13lp32}}: realizar pesquisas e
selecionar informações em fontes diversas, avaliando sua confiabilidade.
\item \textbf{\textsc{em13lp46}}: compartilhar sentidos construídos na leitura
literária, percebendo diferenças entre leituras pessoais e coletivas.
\item \textbf{\textsc{em13lp52}}: analisar obras significativas das literaturas
indígenas, considerando seu contexto de produção, suas matrizes
culturais e o modo como dialogam com o presente.
\item \textbf{\textsc{em13lp54}}: criar obras autorais e produções derivadas
(recontos, ilustrações, entrevistas) em diálogo com o texto literário.
\item \textbf{\textsc{em13lgg401}}: compreender as línguas como fenômeno
(geo)político, histórico, social e cultural.
\item \textbf{\textsc{em13lgg601}}, \textbf{\textsc{em13lgg602}} e \textbf{\textsc{em13lgg603}}:
apropriar-se do patrimônio artístico de diferentes povos, fruí-lo
esteticamente e expressar-se em processos de criação.
\end{itemize}

\textbf{Ciências Humanas e Sociais Aplicadas}

\begin{itemize}
\item \textbf{\textsc{em13chs101}}: comparar diferentes fontes e narrativas, em
diversas linguagens, incluindo tradições orais.
\item \textbf{\textsc{em13chs104}}: analisar a cultura material e imaterial para
identificar conhecimentos, valores e práticas de diferentes sociedades.
\item \textbf{\textsc{em13chs502}}: desnaturalizar e problematizar preconceitos e
discriminações, identificando ações que promovam os direitos humanos.
\item \textbf{\textsc{em13chs601}}: identificar e analisar as demandas e os
protagonismos políticos, sociais e culturais dos povos indígenas no
Brasil contemporâneo.
\end{itemize}

\textbf{Ciências da Natureza e suas Tecnologias}

\begin{itemize}
\item \textbf{\textsc{em13cnt206}}: discutir a importância da conservação da
biodiversidade, relacionando-a aos modos de vida das populações
tradicionais (atividade sobre a floresta do livro).
\end{itemize}

\chapter{Planejamento sugerido}

\vspace*{-3\baselineskip}

A sequência a seguir prevê dez aulas de 50 minutos, além do projeto
integrador, que pode se estender por três ou quatro semanas, em
articulação com outros componentes.

\begingroup\small
\noindent\begin{tabular}{p{.1\linewidth}p{.42\linewidth}p{.38\linewidth}}
\textbf{Aula} & \textbf{Atividades do estudante} & \textbf{Foco}\\
\hline
1 & Antes da leitura: ``O que eu sei, o que eu imagino'' & Levantamento
de conhecimentos prévios e estereótipos\\
2 & Antes da leitura: ``Quem são os Hupd'äh'' e ``Um título, uma imagem,
uma língua'' & Contexto, mapa, língua tonal\\
3 & Durante a leitura: leitura em voz alta e roteiro & Fruição e leitura
atenta\\
4 & Compreensão e interpretação: questões 1 a 7 & Personagens, tempo,
onomatopeias\\
5 & Compreensão e interpretação: questões 8 a 14 & Narrador,
transmissão, narrativa de origem\\
6 & Aprofundamento: oralidade, escrita e tradução; outras gentes &
Autoria coletiva, cosmologia\\
7 & Aprofundamento: línguas; estereótipos & Diversidade linguística,
preconceito\\
8 & Criação e protagonismo: apresentação das propostas e início da
produção & Criação, ética da recriação\\
9 & Pesquisa e debate: mesa-redonda; orientação da redação &
Argumentação oral e escrita\\
10 & Conexões com Goiás; lançamento do projeto integrador & Povos
indígenas de Goiás\\
\end{tabular}
\endgroup

\chapter{Encaminhamentos e respostas esperadas}

\section{antes da leitura}

A lista de palavras (atividade 1) funciona como avaliação diagnóstica.
Registre-a e guarde-a: ela será retomada no bloco ``Ver além dos
estereótipos''. É comum surgirem termos como ``tribo'', ``oca'',
``floresta'', ``cocar'', ``arco e flecha'', ``selvagem''. Não corrija de
imediato; conduza a discussão com as perguntas propostas, levando a turma
a perceber a generalização (``todos iguais''), a cristalização no passado
e a associação exclusiva com a floresta.

Nas questões sobre a nota, espera-se que os estudantes localizem as
informações (Alto Rio Negro, fronteira com a Colômbia; cerca de 1\,500
pessoas em 35 aldeias; caça, pesca, coleta e roças de maniva, com as
mulheres responsáveis pelas roças e participando também da caça, da pesca
e da coleta). Sobre a ``gente'' não humana, acolha hipóteses; o objetivo
é criar expectativa. Na atividade do mapa, o site Terras Indígenas no
Brasil permite localizar as terras do Alto Rio Negro, no município de
São Gabriel da Cachoeira.

Sobre a língua tonal: nessas línguas, a altura melódica de uma sílaba
distingue palavras, isto é, duas palavras com os mesmos sons podem ter
sentidos diferentes conforme o tom. Em português, a entonação atua sobre
a frase, não sobre o significado da palavra.

\section{durante a leitura}

Recomenda-se uma primeira leitura em voz alta feita pelo professor e uma
segunda leitura compartilhada pelos estudantes. Estimule a tentativa de
ler também os trechos em hup com apoio de ``Para ler as palavras hup'':
não se busca a pronúncia correta, mas a experiência de habitar, por um
momento, outra língua.

\section{compreensão e interpretação}

\begin{enumerate}[label=\arabic*.]
\item Espera-se que o estudante identifique a conjunção
\textit{mas}, que introduz uma oposição: a gente-sombra é perigosa, mas
também sábia. Essa ambivalência explica a tensão da história: o homem
hup se arrisca ao se aproximar de um ser perigoso porque ele detém um
saber valioso, os cantos.

\item Espera-se que o estudante perceba que o nome próprio individualiza
e dá importância a Way Naku, que é um ser poderoso e conhecido pelos
Hupd'äh (o canto de abertura leva seu nome). O homem hup, anônimo, pode
representar qualquer Hupd'äh, o que reforça o caráter coletivo da
herança: ele é ``o primeiro'' de todos que depois cantarão.

\item Espera-se que o estudante reconheça que a beleza da música vence o
medo. \textit{Encantado} indica fascínio, atração irresistível, quase
um feitiço; a palavra também remete, no português do Brasil, ao universo
dos seres ``encantados''.

\item O ``Canto da cutia'', o ``Canto grande'', o ``Canto do Umari'' e o
``Canto da gente-sombra''. Espera-se que o estudante observe que os
nomes remetem a animais, frutas e seres da floresta, mostrando a relação
estreita entre a vida cotidiana na mata e a música das festas.

\item a) \textit{Prender} tem aqui o sentido de ``fixar na memória'',
``apreender''; pode ser substituído por \textit{memorizar},
\textit{aprender}, \textit{guardar}. O verbo original sugere algo que se
captura e se retém com esforço. b) Espera-se que o estudante perceba que
a aprendizagem se dá pela escuta prolongada e atenta, pela paciência e
pela memória, sem escrita, ao longo de uma noite e um dia inteiros.

\item Aceite hipóteses coerentes com o texto, por exemplo: o homem hup
quer ir embora sem ser seguido e assusta o homem-sombra para ter tempo
de fugir; quer derrubá-lo para se proteger, pois sabe que a gente-sombra
é perigosa; quer testar ou afastar o ser depois de ter obtido o que
desejava. Valorize a justificativa com base no texto. Pode-se discutir a
ambiguidade ética do gesto: o conhecimento é obtido por escuta e, ao
final, por astúcia.

\item a) Espera-se que o estudante perceba que cada língua representa os
sons segundo seu próprio sistema sonoro e suas convenções: não há uma
única forma ``natural'' de imitar um ruído. b) Porque a onomatopeia tem
função expressiva e sonora; o tradutor precisa escolher entre
transcrever o som original (que pode soar estranho ao leitor) ou recriar
um efeito equivalente na língua de chegada, como fez a tradução
portuguesa.

\item O jacaré é decisivo para o desfecho: ao ser confundido com o
homem hup e morder o braço do homem-sombra, ocupa Way Naku e permite a
fuga do homem hup. Sem ele, o homem-sombra poderia alcançar o homem hup,
e os cantos talvez não chegassem à aldeia.

\item Espera-se que o estudante conclua que Way Naku não enxerga o mundo
como os humanos: confunde um jacaré com uma pessoa. Isso sugere que as
outras gentes percebem os seres de modo diferente e reforça sua
estranheza, sem negar seu poder.

\item Porque explica a origem de uma prática que existe até hoje: os
\textit{caapivaiá}, os cantos das festas hupd'äh. A narrativa liga um
acontecimento do passado a um costume do presente.

\item a) Do homem-sombra ao homem hup, que escuta; do homem hup a seus
filhos, irmãos e cunhados; destes às gerações seguintes, até os Hupd'äh
de hoje, que cantam e dançam nas festas. b) Espera-se que o estudante
reconheça que a frase final legitima a narrativa pela cadeia de
transmissão: o narrador não se apresenta como inventor, mas como elo de
uma memória coletiva; a autoridade da história vem dos avós.

\item A mudança revela que o narrador é um Hupd'äh que fala de dentro da
própria comunidade: a história é dele e do seu povo. A primeira pessoa
aproxima passado e presente e transforma um relato sobre ``um homem
hup'' em memória coletiva (``nossas festas'').

\item No canto, a voz é a do homem-sombra, que colhe cachos de cucura e
os joga para baixo, a mesma ação narrada na história (``ia de galho em
galho cantando e pegando os cachos da fruta''). Os versos finais foram
mantidos em hup porque têm nomes próprios (Way Naku, Marika) e, talvez,
palavras cuja sonoridade e sentido se perderiam na tradução; mantê-los
preserva a musicalidade e lembra ao leitor a língua original.

\item Resposta pessoal. Valorize a comparação entre as hipóteses e o
texto lido e o registro de surpresas, estranhamentos e dúvidas.
\end{enumerate}

\section{aprofundamento e reflexão}

\textbf{Oralidade, escrita e tradução.} Na questão 1,
espera-se a identificação de Mário Andrade Pires (narrador), Ponciano
Salustiano Ramos (cantor do \textit{Way Naku}), os professores hupd'äh
(transcrição e tradução), Patience Epps e Danilo Paiva Ramos
(organização), Anita Ekman (ilustração), além das gerações anteriores
que transmitiram a história. Na questão 2, destaque o direito à educação
escolar indígena diferenciada e bilíngue (Constituição, art. 210, § 2º)
e o fortalecimento da língua e da identidade. Para aprofundar a
discussão sobre oralidade e escrita, é útil o livro \textit{Oralidade e
cultura escrita}, de Walter Ong.

\textbf{Outras gentes.} Estudos etnográficos sobre os Hupd'äh
descrevem um cosmos em camadas (mundo subterrâneo das sombras, a floresta
em que vivemos e o mundo da luz, acima do céu), em que luz e sombra
compõem todos os seres em proporções diversas. Use essa informação para
ampliar a leitura, sem transformá-la em chave única de interpretação.
Na questão sobre seres de matas e rios, conduza a turma a perceber
funções comuns dessas histórias: orientar condutas, alertar para
perigos, regular o uso do ambiente, explicar o desconhecido e reforçar
laços comunitários.

\textbf{Muitas línguas.} A cooficialização garante o uso da
língua em serviços públicos, documentos e campanhas, a formação de
servidores falantes e o ensino. Além de São Gabriel da Cachoeira, outros
municípios cooficializaram línguas indígenas nas últimas décadas;
oriente a pesquisa em fontes confiáveis e peça que os estudantes
indiquem as leis municipais correspondentes.

\textbf{Estereótipos.} A frase sobre o celular é
preconceituosa porque associa a identidade indígena a uma imagem
congelada no passado e nega aos povos indígenas o direito de usar
tecnologias, como qualquer outro grupo. A identidade étnica se baseia no
autorreconhecimento e no pertencimento a um povo, não em aparência ou
em objetos. A última questão permite discutir que tradição não é
imobilidade: registrar histórias em livros, com parceiros, é também uma
estratégia contemporânea de continuidade cultural.

\section{criação e protagonismo}

Apresente todas as propostas e permita que cada estudante escolha pelo
menos uma, garantindo que a turma, no conjunto, contemple as quatro
primeiras (ilustração, reconto, entrevista, pesquisa), como prevê o
Termo de Referência. Insista nos cuidados éticos: citar a fonte, não
copiar grafismos ou cantos como enfeite e evitar caricaturas. Na roda de
cantos e histórias, oriente os estudantes a pedir autorização antes de
gravar e a identificar quem contou cada história.

\section{pesquisa e debate}

Na mesa-redonda, espera-se que surjam argumentos como: a escrita
preserva e permite a circulação e a alfabetização na língua; mas fixa
uma versão, retira a voz e o contexto, e pode deslocar a autoridade dos
anciãos para quem escreve. O ponto central é a participação e o controle
das comunidades sobre o registro e a circulação de seus saberes. Na
redação, oriente a construção de um projeto de texto com tese clara,
dois argumentos desenvolvidos (por exemplo, a transmissão
intergeracional ameaçada e a falta de políticas públicas de ensino
bilíngue) e proposta de intervenção completa (agente, ação, meio,
finalidade e detalhamento).

\section{conexões com goiás}

Peça que os estudantes registrem as fontes. O site Terras Indígenas no
Brasil e a enciclopédia Povos Indígenas no Brasil trazem verbetes sobre
os Karajá, os Avá-Canoeiro e os Tapuia. O dossiê do Iphan sobre as
bonecas \textit{ritxoko} está disponível no portal do instituto. Na
atividade de toponímia, alerte para as etimologias populares sem base:
peça que os estudantes comparem fontes e indiquem quando a origem é
incerta.

\chapter{Cuidados na mediação}

\section{laicidade e respeito}

A narrativa trata de seres, poderes e cantos que fazem parte da
cosmologia hupd'äh. Na escola pública e laica, o objetivo não é aderir
nem converter, mas compreender: apresentar essas narrativas como
sistemas de conhecimento legítimos, em pé de igualdade com as narrativas
de outras tradições culturais e religiosas presentes na turma. Evite
tanto a exotização (tratar como curiosidade estranha) quanto a
desqualificação (``crendice'', ``superstição''). Se algum estudante
manifestar desconforto por motivos religiosos, acolha e explique o
caráter de estudo cultural e literário da atividade.

\section{termos a evitar ou ressignificar}

\begin{itemize}
\item \textit{Índio}, de modo genérico: prefira ``indígena'', ``povos
indígenas'' ou, sempre que possível, o nome do povo (Hupd'äh, Karajá,
Guarani Mbya).
\item \textit{Tribo}: prefira ``povo'' ou ``comunidade''.
\item \textit{Lenda}, \textit{folclore}, \textit{crendice}: prefira
``narrativa'', ``narrativa de origem'', ``cosmologia''.
\item \textit{Primitivo}, \textit{selvagem}: carregam a ideia falsa de
estágio inferior de desenvolvimento.
\item \textit{Descobrimento do Brasil}: apaga a presença milenar dos
povos originários; prefira ``chegada dos portugueses'', ``invasão'',
``início da colonização''.
\item \textit{Cara de índio}, \textit{cabelo de índio}: generalizam
traços físicos e ignoram a diversidade dos povos indígenas.
\item \textit{Maku}: nome externo e pejorativo; use Hupd'äh e família
naduhup.
\end{itemize}

\section{diante de falas preconceituosas ou bullying}

Comentários preconceituosos podem surgir, às vezes em tom de brincadeira.
Não os ignore: interrompa com serenidade, nomeie o problema (``essa fala
reproduz um estereótipo''), devolva a pergunta à turma e retome as
informações do livro e das pesquisas. Se houver estudantes indígenas ou
de famílias indígenas na turma, não os exponha nem os coloque na posição
de ``representantes'' do tema; se desejarem compartilhar experiências,
acolha. Casos de discriminação ou bullying de caráter étnico-racial
devem ser registrados e encaminhados à coordenação pedagógica, conforme
os protocolos da unidade escolar, com envolvimento das famílias.

\chapter{Avaliação}

\vspace*{-4\baselineskip}

A avaliação deve ser processual e considerar a participação nas leituras
e discussões, o diário de leitura, as respostas às questões, as
produções criativas, o debate, a redação e o projeto integrador. A
retomada da lista de palavras feita antes da leitura permite observar
mudanças de percepção e pode compor a autoavaliação.


\begingroup\small
\noindent\begin{tabular}{p{.22\linewidth}p{.23\linewidth}p{.23\linewidth}p{.2\linewidth}}
\textbf{Critério} & \textbf{Pleno} & \textbf{Adequado} & \textbf{Em
desenvolvimento}\\
\hline
Compreensão da obra & Dialoga com a narrativa e seu contexto de modo
preciso e criativo & Retoma a narrativa corretamente, com pouca
ampliação & Apresenta equívocos sobre a narrativa ou o povo\\
\hline
Uso de fontes & Cita o livro e fontes pesquisadas, comparando-as &
Cita o livro e ao menos uma fonte & Não identifica as fontes\\
\hline
Ética e respeito & Respeita autoria e significados; evita estereótipos
& Respeita a autoria, com algum estereótipo pontual & Reproduz
estereótipos ou copia elementos sem contexto\\
\hline
Linguagem e forma & Recursos expressivos adequados ao gênero escolhido
& Linguagem adequada, com recursos pouco explorados & Linguagem
inadequada ao gênero ou à situação\\
\hline
Colaboração & Contribui ativamente e apoia o grupo & Cumpre sua parte
& Participa pouco das etapas\\
\end{tabular}
\endgroup

\chapter{Referências}

\begin{bibliohedra}
\tit{brasil}. \textit{Constituição da República Federativa do
Brasil de 1988}. Brasília: Senado Federal, 1988.
\tit{brasil}. Lei nº 11.645, de 10 de março de 2008. Altera a
Lei nº 9.394/1996, modificada pela Lei nº 10.639/2003, para incluir no
currículo oficial da rede de ensino a obrigatoriedade da temática
``História e Cultura Afro-Brasileira e Indígena''.
\tit{brasil}. Ministério da Educação. \textit{Base Nacional
Comum Curricular}. Brasília: MEC, 2018.
\tit{epps}, Patience. \textit{A Grammar of Hup}. Berlim: Mouton
de Gruyter, 2008.
\tit{goiás}. Secretaria de Estado da Educação.
\textit{Documento Curricular para Goiás -- Etapa Ensino Médio}.
Goiânia: Seduc-GO, 2021.
\tit{graúna}, Graça. \textit{Contrapontos da literatura
indígena contemporânea no Brasil}. Belo Horizonte: Mazza, 2013.
\tit{instituto socioambiental}. Povos Indígenas no Brasil:
Hupda. Disponível em: \textit{https://pib.socioambiental.org}.
\tit{krenak}, Ailton. \textit{Ideias para adiar o fim do
mundo}. São Paulo: Companhia das Letras, 2019.
\tit{ong}, Walter. \textit{Oralidade e cultura escrita: a
tecnologização da palavra}. Campinas: Papirus, 1998.
\tit{ramirez}, Henri. \textit{A língua dos Hupd'äh do Alto Rio
Negro}. São Paulo: Associação Saúde Sem Limites, 2006.
\tit{ramos}, Danilo Paiva. \textit{Círculos de coca e fumaça:
encontros noturnos e caminhos vividos pelos Hupd'äh}. São Paulo: Hedra.
\tit{viveiros de castro}, Eduardo. Os pronomes cosmológicos e o
perspectivismo ameríndio. \textit{Mana}, Rio de Janeiro, v.~2, n.~2,
p.~115--144, 1996.
\tit{wittmann}, Luisa Tombini (org.). \textit{Ensino (d)e
história indígena}. Belo Horizonte: Autêntica, 2015.
\end{bibliohedra}
